% ======================================================================
% Section 7 — The higher rungs: T=7 and T=8
% ======================================================================

\section{The higher rungs: $T=7$ and $T=8$}\label{sec:higherrung}

The closures of \S\ref{sec:rigid}--\S\ref{sec:pq} live at the rung $T=6$
($n=5$): four bad sets per pair, the frontier cell $(1K,3U)$, and the
fiber problem is the three-AP census classified by Conjecture~\ref{conj:SL}.
Open Problem~9 asked whether the method survives a third rung. This
section answers affirmatively, with measurements: the machinery, the
classification, and the shear mechanism all port, and at the \emph{critical}
cell of $T=8$ the census is sharper than at $T=6$ --- exactly eleven
families with placement counts \emph{constant in $p$}. The higher rungs
are also where the pinned sets grow, and the shear step must be
re-established as a growth comparison; we measure that growth and the
drift-exit mechanism that keeps it harmless.

\subsection{The rung map and the general-\texorpdfstring{$T$}{T} fiber machinery}

At rung $T=n+1$ the composite modulus is $N=p^2$ with $c=\floor{(N-1)/T}$
and bad sets $B_w=\{x:\wn{wx}\le c\}$. Writing $p=Tk+\epsilon$
($\epsilon=p\bmod T$), a kernel member $pa$ covers exactly the fibers of
$a^{-1}B_k$, where $B_k=\{r:\min(r,p-r)\le k\}$ has $2k+1$ elements, so
the uncovered set has $|U|=(T-2)k+\epsilon-1$. Cells are families of
$n-1=T-2$ bad sets; the frontier cells are
\[
T=7:\ (1K,4U)\ \text{(overlap }k+O(1)\text{)},\qquad
T=8:\ (2K,4U)\ \text{(overlap }O(1)\text{)}\ \text{and}\ (1K,5U),
\]
$(2K,3U)$ at $T=7$ and $(3K,3U)$ at $T=8$ being dead by size
($3(2k+2)<p$ for $k$ past a small boundary). The $T=8$ cell $(2K,4U)$ is
\emph{critical} in the $T=6$ sense --- four footprints of size
$2k+1\approx p/4$ sum to $p+O(1)$ --- which is why its census mirrors
(and sharpens) the base rung's.

The fiber machinery is $T$-free in form; we record it once.

\begin{lemma}[footprints, the three-tier arc, and the drift at rung $T$]
\label{lem:rungT}
Fix rung $T$, $p=Tk+\epsilon$ prime, and a unit $u$ with residue
$r=u\bmod p$ and lift $m$ ($u=r+pm$). On the fiber $F_j=\{j+pt\}$, writing
$uj=a_1p+a_0$, the trace $B_u\cap F_j$ in the $t$-coordinate is an AP
with difference $\bar u=u^{-1}\bmod p$, size
\[
L(a_0)\in\{2k,\ 2k+1,\ 2k+2\}
\]
given by the \emph{three-tier arc law}: with $M=k\epsilon+\delta_0$,
$\delta_0=\floor{(\epsilon^2-1)/T}$, one has $L=2k+1$ on the ball $B_k$
\emph{and} on two shallow off-bands, $L=2k+2$ on the deep band
$a_0\in(p-1-M,\,M]$ when $2M\ge p-1$, and $L=2k$ on the deep band
$a_0\in(M,\,p-1-M]$ when $2M\le p-2$. At $\epsilon\equiv\pm1\pmod T$ the
deep band is the whole off-ball (the phenomenon of \S\ref{sec:rigid});
at intermediate $\epsilon$ a shallow band of width $|M-k|$ carries the
on-size. Consequently the occurring footprint sizes are exactly two
values, $\{2k,2k+1\}$ for $\epsilon<T/2$ and $\{2k+1,2k+2\}$ for
$\epsilon>T/2$. The starts obey the \emph{drift law}
\[
\sigma_u(j)\;=\;\sigma_r^0(j)-\bar u\,m\,j\pmod p,
\qquad \sigma_r^0(j)=\bar u\,(S_2(a_0)-\floor{rj/p}),
\]
so the lift acts by per-fiber $t$-translation linear in $j$; for a
$\pm$-colliding pair $w=u+pm$ the relative drift
$\sigma_w(j)-\sigma_u(j)=-\bar u\,m\,j$ is exactly linear, injective in
$j$, and vanishes only when $B_w=B_u$.
\end{lemma}

\begin{proof}
The footprint computation is the one of \S\ref{sec:composites} verbatim
with $6$ replaced by $T$: $x\in B_u$ iff $a_0+p\,s$ (with
$s=(a_1+ut)\bmod p$) lands in $[-c,c]$ modulo $p^2$, and the good $s$
form the cyclic interval $[S_2,p-1]\cup[0,S_1]$ with
$S_1=k+\floor{(M-a_0)/p}$, $S_2=p-k-\floor{(M+a_0)/p}$, which gives the
three-tier law by the case split on $a_0$ relative to $M$ and $p-1-M$.
The drift law follows from $\floor{(r+pm)j/p}=\floor{rj/p}+mj$ and the
$a_0$-dependence of the arc being $r$-only. Machine-verified: the law
reproduces the committed $T=6$ tables boundary for boundary
($p\le37$, both classes); it holds for every $a_0$ at $19$ primes of
$T=7/T=8$ (all residue classes); the footprint formula matches direct
$Z_{p^2}$ enumeration ($957/957$); the drift identity holds in
$3{,}528/3{,}528$ checks at $T\in\{7,8\}$.
\end{proof}

\subsection{The critical census at \texorpdfstring{$T=8$}{T=8}: eleven families, constant counts}

The $(2K,4U)$ fiber problem --- four APs with sizes in the two-value
window covering $\Z_p$, normal form $(0,1,s_1)$ plus three moving APs
with canonical differences --- is the exact analogue of the
Conjecture~\ref{conj:SL} census one rung up, and it is \emph{finite}.

\begin{proposition}[the 4-difference partner representatives; machine]
\label{prop:T8census}
At every prime $p=8k+\epsilon$ tested --- seventeen primes
$p\in\{19,29,31,41,43,47,59,61,67,71,73,79,89,97,103,137\}$ plus the
boundary $p=17$, covering all four odd $\epsilon$ classes with $k$ up to
$17$ --- the census of four-AP coverings in the normal form returns
\emph{exactly eleven difference families}:
\[
(1,1,1),\quad\text{the permutations of }(1,1,t),\ (1,2,2),\ (1,q,q),
\quad(3,3,3),
\]
where $q=(p-1)/2$ is the antipode and $t=\floor{(p+1)/3}$ (equivalently
$t\equiv\pm3^{-1}\bmod p$, $3t\in\{p-1,p+1\}$). The placement counts are
independent of $p$ and depend only on the class and the $\epsilon$-class
of $p$:
\begin{center}
\small
\begin{tabular}{@{}lccc@{}}
\toprule
$\epsilon$ class & families & placements per family & total \\
\midrule
$\epsilon\in\{3,7\}$ & $(1,1,1)$: $48$; each other: $24$ & $288$ \\
$\epsilon\in\{1,5\}$ & three classes: $528$, $232$, $216$ & $2{,}784$ \\
\bottomrule
\end{tabular}
\end{center}
\noindent No pairwise $\pm$-distinct covering occurs at any $k\ge2$
($\epsilon\in\{3,5,7\}$) or $k\ge5$ ($\epsilon=1$); the $k=2$ boundary
carries twelve $\pm$-distinct configurations, all enumerated. At
$\epsilon=5$, $k=3$ only, twelve further families with six placements
each occur; they are gone by $k=7$.
\end{proposition}

The five mechanisms are hand-derived, and each count is the junction
arithmetic of a placement-rigid structure --- the Corollary~T1
phenomenon two rungs up. \textbf{(M1)} \emph{Four intervals} $(1,1,1)$:
the partition multiset $\{2k+1,(2k+2)^3\}$ gives $4$ ordered size
tuples $\times\,3!=6$ cyclic arrangements $=24$ exact tilings; the
one-overlap multiset $(2k+2)^4$ gives $3!\times4=24$; total $48$.
\textbf{(M2)} \emph{Two intervals + the parity split} $(1,2,2)$: the
intervals tile an arc of length $4k+3$; the complementary arc of length
$4k+4$ splits into its two parity classes, each \emph{being} a 2-AP of
size $2k+2$ --- the $(2,2)$ mechanism of \S\ref{sec:gt-census}, one rung
up, with an interval pair in place of the single interval.
\textbf{(M3)} \emph{Three intervals + the three-block sampler}
$(1,1,t)$: since $t\equiv\pm3^{-1}$, the $t$-AP decomposes into three
interleaved strands --- three short blocks spaced $\approx p/3$ apart
--- the wrapped two-block one block up; the three intervals tile the
complement of the three holes. This is why the \emph{third-point}
replaces the antipode in the double-interval family: three intervals
leave three holes. \textbf{(M4)} \emph{Intervals + antipodal
two-blocks} $(1,q,q)$: the $q$-AP is the wrapped two-block of
\S\ref{sec:gt-census} verbatim; the interval plugs its gap structure.
\textbf{(M5)} \emph{The thickened arc} $(3,3,3)$: three 3-APs with
starts spaced $1,2$ have union a contiguous arc of length $3s-2=6k+4$;
the normalized interval covers the complementary arc of length $2k+3$.
Every count is $p$-free because each configuration is determined by its
junction pattern.

\subsection{The \texorpdfstring{$T=7$}{T=7} census: the core, the tail, and the growing pinned sets}

At $T=7$ the frontier cell $(1K,4U)$ carries structural linear slack
(the overlap window is $[k-\epsilon+4,\,k-\epsilon+8]$), and the census
is a \emph{core--tail} object rather than a finite list.

\begin{proposition}[the $T=7$ census; machine]\label{prop:T7census}
At eighteen primes $p=7k+\epsilon$, all six $\epsilon$ classes, $k$ up
to $11$: (i) \emph{no pairwise $\pm$-distinct four-AP covering occurs}
at any $k\ge4$ ($\epsilon\le3$) resp.\ $k\ge2$ ($\epsilon\ge4$) --- the
clause-(i) analogue of Conjecture~\ref{conj:SL} one rung up; the
boundary dirt at $k\le3$ is finite and enumerated. (ii) The same eleven
families as Proposition~\ref{prop:T8census} --- with
$t=\floor{(p+1)/3}$, $q=(p-1)/2$ --- are the top eleven by
configuration count at every prime and carry $\approx90\%$ of all
configurations at $k\ge6$. (iii) The core's placement counts grow with
the overlap window (the $(1,1,1)$ count is
$528,2160,5712,11952,21648,35568$ at $k=2,\dots,11$: quadratic growth,
the junction freedom under the slack; the count depends only on the
window, whence exact coincidences between primes with equal windows).
(iv) Beyond the core, $30$--$44$ tail families occur at $k\ge5$ (down
from $120$--$170$ at $k\le4$), with rational-anchor differences
($\{4,5,6,12,14,18,28,33\}$ at $p=71$) and counts a factor $15$--$50$
below the core.
\end{proposition}

This is the measured form of the ``growing pinned sets'': the solution
sets grow $\approx k^2$ (two-dimensional families), and the $T=6$
pigeonhole $|\Delta|\le6<|U|$ must be replaced by a growth comparison.
The replacement is supplied by the projections and the exit time.

\subsection{The shear step at the higher rungs}

\begin{proposition}[projections, exit, and the empty cells; machine]
\label{prop:shearT78}
\textup{(i)} The one-coordinate projections of the census solution sets
(starts and relative offsets) have maximum size
\[
11,\,12,\,18,\,24,\,24,\,30,\,36,\,36,\,42,\,42,\,42
\quad\text{at }k=1,\dots,11,
\]
that is, \emph{saturating at $42$}, while
$|U|=5k+\epsilon-1$ grows linearly ($20\to60$ over the same range); the
pigeonhole $|U|>\text{maxproj}$ holds at every $k\ge5$ except the single
case $(p{=}29,\epsilon{=}1,k{=}4)$, where $20<24$. \textup{(ii)} The
exit time of an affine drift line from a census solution family is at
most $2$ consecutive fibers, over all drift vectors fixed by the first
fiber (top-ten families at $p\in\{17,29,43\}$; all admissible families
in the ground-truth runs). \textup{(iii)} Ground truth, exhaustive over
census-admissible families with the bad-set distinctness filter: the
$(1K,4U)$ cell at $T=7$ is \emph{empty} at every prime
$11\le p\le61$ (all six $\epsilon$ classes; best exit $2$); the
$(2K,4U)$ cell at $T=8$ is \emph{empty} at every prime $17\le p\le43$
tested (best exit $\le1$; at $p=29,41$ no admissible family even passes
the per-fiber size profile). Controls: brute force at $p\in\{5,7,11\}$
(all unit triples, distinct bad sets, no census restriction) finds
$262{,}350$ triples and $0$ coverings at $p=11$, and $8$ coverings at
the degenerate $p=5$ ($k=0$) --- exactly the $8$ found by the
restricted search; $0$ at the $\epsilon=0$ prime $p=7$.
\end{proposition}

The closure architecture is the $T=6$ one, rung-shifted: the census
forces a $\pm$-collision (Propositions~\ref{prop:T8census}--\ref{prop:T7census});
the colliding pair's relative offset drifts exactly linearly and
injectively (Lemma~\ref{lem:rungT}); the relative-offset projections are
bounded --- measured at $\le42$, with the bound's mechanism the
span-exact junction arithmetic of M1--M5, the same analysis that proves
Corollary~T1 --- and $|U|$ is linear in $k$; the drift line exits the
solution family within two fibers. The degenerate $p=5$ ($k=0$) is the
boundary phenomenon, recorded as at $T=6$.

The $(1K,5U)$ cell named below is the object of the conditional
closing attack of \S\ref{sec:scale} (Theorem~\ref{thm:SC}); the honest
boundary recorded here is the measured boundary there.

\subsection{The honest boundary of the port}

Three things do not port, and we record them. \emph{First}, the
$(1K,5U)$ cell at $T=8$: the five-difference census at $k=2$ has $8{,}168$
families and $32.5$M configurations, $840$ of them $\pm$-distinct --- the
collision-only law fails at the fifth difference, the census restriction
does not confine this cell, and its closure (sampled: best run $2$ over
$20{,}000$ random families at $p\in\{17,19\}$; heuristic: solution
families of dimension $3$ in a $4$-coordinate relative space still meet a
generic affine line boundedly) is the rung's open piece. \emph{Second},
the $T=7$ tail: the family count decays ($183\to43$) but is not proved
to die; only the $T=8$-critical census is a finite exact list.
\emph{Third}, the pinning lemma --- projections $\le42$ --- is measured
with the mechanisms identified but not yet proved; the span-exactness
proof per family is the recorded obligation, the direct analogue of the
$T=6$ pinning-size lemma (Input~S's fiberwise form).

For the programme's single unbounded input the meaning is this: the
four-family classification of Conjecture~\ref{conj:SL} is not a $T=6$
accident. One rung up the same anchors occur ($1$, $2$, $q$) with a new
pair ($3$, $t$); the counts stay $p$-independent; the exhaustiveness
clause survives; and the shear mechanism that consumes the
classification is $T$-free in its linear part. The two-page computation
that would discharge Input~S now has its higher-rung shape: prove the
eleven-family exhaustiveness at $T=8$-critical first (a finite, exactly
verified statement), then the projection bound, and the $T=6$ statement
--- the base rung of the same phenomenon --- follows by the same
mechanism analysis at one lower rung.
